\chapter{Exigences fonctionnelles}

\begin{longtable}{>{\centering\arraybackslash}p{2cm} >{\raggedright\arraybackslash}p{12cm}}
    \caption{Exigences fonctionnelles du système}\label{tab:ef}\\
    \toprule
    \rowcolor{violetprincipal}
    \thead{Référence} & \thead{Exigence} \\
    \midrule
    \endfirsthead
    \toprule
    \rowcolor{violetprincipal}
    \thead{Référence} & \thead{Exigence} \\
    \midrule
    \endhead
    \bottomrule
    \endfoot
    EF01 & Créer un document textuel. \\
    \rowcolor{violettresclair}
    EF02 & Partager un document avec d'autres rédacteurs. \\
    EF03 & Édition simultanée d'un document. \\
    \rowcolor{violettresclair}
    EF04 & Propager les modifications en quasi-temps réel. \\
    EF05 & Indiquer les participants présents et leurs positions d'écriture. \\
    \rowcolor{violettresclair}
    EF06 & Conserver l'historique des versions. \\
    EF07 & Consulter et restaurer les versions antérieures. \\
    \rowcolor{violettresclair}
    EF08 & Créer un lien de téléchargement avec un quota défini. \\
    EF09 & Produire un PDF correspondant à l'état retenu du document au moment de la prise en charge. \\
    \rowcolor{violettresclair}
    EF10 & Garantir le respect strict du quota de téléchargement. \\
    EF11 & Appliquer le quota quotidien d'écriture de chaque rédacteur. \\
    \rowcolor{violettresclair}
    EF12 & Détecter et traiter les tentatives de contournement des quotas. \\
    EF13 & Permettre à l'association de superviser l'état de l'application. \\
    \rowcolor{violettresclair}
    EF14 & Notifier le créateur d'un lien lorsque sa dernière utilisation autorisée est consommée. \\
    EF15 & Informer les visiteurs de l'état de leur demande de téléchargement. \\
\end{longtable}

